\documentclass[10pt]{amsart}
\usepackage[a4paper,margin=23mm]{geometry}
\usepackage{amsmath,amssymb,amsthm,mathtools,hyperref,graphicx,xspace,letltxmacro,xpatch}
\usepackage[english]{babel}
\newtheorem{theo}{Theorem}[section]
\newtheorem{lemm}[theo]{Lemma}\newtheorem{prop}[theo]{Proposition}\newtheorem{coro}[theo]{Corollary}\newtheorem{conj}[theo]{Conjecture}
\theoremstyle{definition}\newtheorem{defi}[theo]{Definition}\newtheorem{rema}[theo]{Remark}\newtheorem{exam}[theo]{Example}\newtheorem{exem}[theo]{Example}\newtheorem{nota}[theo]{Notation}
\newcommand{\ie}{i.e.}\newcommand{\eg}{e.g.}\newcommand{\resp}{resp.}\newcommand{\skpt}{\par\smallskip\noindent}
\emergencystretch=3em\pagestyle{plain}
\usepackage{subfig}
\usepackage{xpatch}

\DeclareMathOperator{\E}{\mathbf{E}}



\newcommand{\bbN}{\mathbb{N}}
\newcommand{\Z}{\mathbb{Z}}
%\newcommand{\R}{\mathbb{R}}
%\newcommand{\C}{\mathbb{C}}
\newcommand{\T}{\mathbb{T}}
\newcommand{\ind}[1]{\mathbf{1}_{\left\{#1\right\}}}
%\newcommand{\indset}[1]{\mathbf{1}_{#1}}
\newcommand{\floor}[1]{{\left\lfloor #1 \right\rfloor}}
%\newcommand{\ceil}[1]{{\left\lceil #1 \right\rceil}}

\renewcommand{\P}{\mathbf{P}}


\newcommand{\defeq}{:=}

\renewcommand{\S}{\mathcal{S}}
\newcommand{\clH}{\mathcal{H}}
\newcommand{\clF}{\mathcal{F}}
\newcommand{\clG}{\mathcal{G}}
%%%%%%%%%%%%%%%%%%%%%%%%%%%%%%%%%%%%%%%%%%%%%%%%%%%%%%%%%%%
%~ This fixes spacing issues around \left( and \right) and other similar symbols
\let\originalleft\left 
\let\originalright\right
\renewcommand{\left}{\mathopen{}\mathclose\bgroup\originalleft}
\renewcommand{\right}{\aftergroup\egroup\originalright}
%%%%%%%%%%%%%%%%%%%%%%%%%%%%%%
\let\oldin\in
\renewcommand{\in}{\mathchoice{\oldin}{\oldin}{\,\oldin\,}{\,\oldin\,}}
%%%%%%%%%%%%%%%%%Arrow
\renewcommand*{\to}{\mathchoice{\longrightarrow}{\rightarrow}{\rightarrow}{\rightarrow}}
\let\oldmapsto\mapsto
\renewcommand*{\mapsto}{\mathchoice{\longmapsto}{\oldmapsto}{\oldmapsto}{\oldmapsto}}

%%%%%%%%%%%%%%%%%%%%%%%%%%%%%%%%%%%%%%%%%%%%%%%%%%%%%%%%%%%%%%%%%%%%
\graphicspath{{./figures/}}

\newcommand*{\mk}{\mkern -1mu}
\newcommand*{\Mk}{\mkern -2mu}
\newcommand*{\mK}{\mkern 1mu}
\newcommand*{\MK}{\mkern 2mu}

%\hypersetup{urlcolor=purple, linkcolor=blue, citecolor=red}

\newcommand*{\relabel}{\renewcommand{\labelenumi}{(\theenumi)}}
\newcommand*{\romanenumi}{\renewcommand*{\theenumi}{\roman{enumi}}\relabel}
\newcommand*{\Romanenumi}{\renewcommand*{\theenumi}{\Roman{enumi}}\relabel}
\newcommand*{\alphenumi}{\renewcommand*{\theenumi}{\alph{enumi}}\relabel}
\newcommand*{\Alphenumi}{\renewcommand*{\theenumi}{\Alph{enumi}}\relabel}
\let\oldtilde\tilde
\renewcommand*{\tilde}[1]{\mathchoice{\widetilde{#1}}{\widetilde{#1}}{\oldtilde{#1}}{\oldtilde{#1}}}
\let\oldhat\hat
\renewcommand*{\hat}[1]{\mathchoice{\widehat{#1}}{\widehat{#1}}{\oldhat{#1}}{\oldhat{#1}}}

\newcommand{\llbracket}{\mathopen{[\mkern -2.7mu[}}
\newcommand{\rrbracket}{\mathclose{]\mkern -2.7mu]}}
\newcommand{\llbracketclose}{\mathclose{[\mkern -2.7mu[}}
\newcommand{\rrbracketopen}{\mathopen{]\mkern -2.7mu]}}

\newcommand{\llb}{\llbracket}
\newcommand{\rrb}{\rrbracket}
\newcommand{\llbigbracket}{\mathopen{\bigl[\mkern -4.2mu\bigl[}}
\newcommand{\rrbigbracket}{\mathclose{\bigr]\mkern -4.2mu\bigr]}}
\newcommand{\llbg}{\llbigbracket}
\newcommand{\rrbg}{\rrbigbracket}
%%%%%%%%%%%%%%%%%%%%%%%%%%%%%%%%%%%%%%%%%%%%%%%%%%%%%%%%%%%%%%%%%%%%

\makeatletter
\expandafter\def\csname externalRef@fig:IBM\endcsname{2.1}
\expandafter\def\csname externalRef@fig:couplage\endcsname{3.1}
\expandafter\def\csname externalRef@thm:finitevinfinite\endcsname{1.9}
\expandafter\def\csname externalRef@theo1.9.1\endcsname{1}
\LetLtxMacro\NativeRef\ref
\newcommand{\SourceRef}[1]{\ifcsname externalRef@#1\endcsname\href{https://ahl.centre-mersenne.org/item/AHL_2025__8__635_0/}{\csname externalRef@#1\endcsname\ of the source article}\else\NativeRef{#1}\fi}
\AtBeginDocument{\LetLtxMacro\NativeRef\ref\LetLtxMacro\ref\SourceRef}
\makeatother
\title{Uniform expected terminal occupancy for every bounded initial infinite-bin configuration}
\author{Bastien Mallein, Sanjay Ramassamy, Arvind Singh}
\date{doi:10.5802/ahl.244:published; CC BY 4.0}
\begin{document}
\maketitle
\noindent\textit{Editorial note.} The stronger uniform expectation proposition and its almost-sure freezing consequence are one counted problem. Figure illustrations are not necessary for the defined rank-based dynamics. Author proof is retained; this focused review does not independently rederive each conditional expectation.\par\medskip
\noindent\textit{Allowed standard prerequisites (editorial).}\begin{itemize}
\item Conditional expectation and countable stochastic-process construction. Ordinary conditioning, Tonelli, Fatou and monotone convergence; the tail contraction and uniform cutoff comparison are not given as free inputs.
\end{itemize}\medskip
We introduce some notation needed to define the infinite-bin model formally. The state space of the IBM is defined as
\[
\S
:= \bigl\{ X : \Z \to \Z_+\cup \{\infty\} \text{ such that } \sup\{ j \in \Z : X(j) \neq 0 \} \in (-\infty,\infty) \bigr\}.
\]
An element $X\in \S$ represents a \emph{configuration of balls} where, for each $k \in \Z$, the bin at position $k$ contains $X(k)$ balls.

Let us point out that the total number of balls in a configuration $X \in \S$ may be finite or infinite (but cannot be zero). We even allow the number of balls in a given bin $k$ to be infinite i.e. $ X(k) = \infty$. For the dynamics of the infinite bin model, such a bin will represent an impassable barrier (the configuration on the left of this bin will not matter and remain unchanged forever). A configuration $X \in \S$ is said to be
\begin{itemize}
\item \emph{locally finite} if $X(k) < \infty$ for all $k\in \Z$,
\item \emph{bounded} if $\sup_{k\in \Z} X(k) < \infty$.
\end{itemize}

Given $\mu$ a probability distribution on $\bbN$ and $X_0 \in \S$ an initial configuration, an infinite-bin model $X$ with reproduction law $\mu$ starting from $X_0$ is constructed as an $\S$-valued Markov process in the following fashion. Given $(\xi_n, n \in \bbN)$ i.i.d. random variables with law $\mu$, at each integer time $n$, $X_n$ is constructed from $X_{n-1}$ by adding a new ball to the bin immediately to the right of the bin containing the $\xi_{n}^{\rm th}$ rightmost ball. In other words, writing $J_{n-1}$ for the index of the bin containing the $\xi_{n}^{\rm th}$ rightmost ball at time $n-1$, defined by
\[
\sum_{j = J_{n-1}}^\infty X_{n-1}(j) \geq \xi_{n} > \sum_{j=J_{n-1} +1}^\infty X_{n-1}(j),
\]
we set\footnote{If such an index does not exist, which is the case if the configuration $X_{n-1}$ contains fewer than $\xi_{n}$ balls, we just set $X_{n}=X_{n-1}$.} $X_{n}(\cdot) = X_{n-1}(\cdot) + \ind{\cdot\,=\,J_{n-1}+1}$.

In contrast, we focus here on a different question and study instead how the number of balls in a given bin grows as time passes. Obviously, for all $k \in \Z$, the number $X_n(k)$ of balls in the bin at position $k$ at time $n$ is non-decreasing in $n$ (since balls are never removed from bins). Hence we can define the random variables
\[
X_\infty(k) = \lim_{n \to \infty}\uparrow X_n(k) \in \Z_+ \cup \{\infty\} \quad \text{a.s.}
\]
We call the event $\{X_\infty(k) < \infty\}$ the \emph{freezing} of bin $k$, as it means that, on this event, there is some finite (random) time after which no ball is ever added to bin $k$. This paper aims to describe conditions on the distribution $\mu$ and the initial configuration $X_0$ for which freezing occurs.

\begin{theo}\label{thm:unifbounded}
Let $X_0 \in \S$ be a bounded configuration. For any distribution~$\mu$, we have $X_\infty(k) < \infty$ a.s. for all $k\in \Z$.
\end{theo}

Given an infinite-bin model configuration $X \in \S$, we define its \emph{rightmost barrier} as
\begin{equation}\label{def:deltabarrier}
\delta(X)
:= \sup\{j\in \Z : X(j) = +\infty\} \;\in \Z\cup\{-\infty\}
\end{equation}
with the convention $\delta(X) = \sup \emptyset = -\infty$ when $X$ is locally finite. The \emph{front} of $X$ is the location of its rightmost non-empty bin (which is well-defined by definition of a configuration), i.e.
\[
F(X)
:= \max\{ j \in \Z : X(j) \neq 0\} \in \Z.
\]
For all $k \in \bbN$, we also define $B(X,k)$ as the location of the $k^{\rm th}$ rightmost ball in $X$, i.e.
\[
B(X,k)
:= \sup\left\{ j \in \Z : \sum_{i=j}^\infty X(i) \geq k \right\},
\]
again with $\sup \emptyset = -\infty$. We observe that
\[
F(X) = B(X,1)\quad\text{and}\quad\delta(X) = \lim_{k\to\infty} B(X,k).
\]

\begin{rema}
For all $X \in \S$ and $k \in \bbN$, the quantity $B(X, k)$ is well-defined irrespectively of the relative ordering of the balls inside each bin. In other words, the dynamic of the infinite-bin model does not require to fix an ordering of the balls. However, it will be convenient to specify an ordering for the balls in each bin, in order to introduce a genealogical structure to the particle system.
\end{rema}


To give a formal definition of the infinite-bin model dynamics, we introduce the operator $\Phi_\xi : \S \to \S$ that maps a configuration to the new configuration obtained by adding a single new ball immediately to the bin on the right of the $k^{\rm th}$ rightmost ball:
\begin{equation}\label{eq:operatorPhi}
\Phi_\xi(X)
:=
\begin{cases}
\left(X(j) + \ind{j\,=\,B(X,\xi)+1}, j \in \Z\right) & \text{ if } B(X,\xi) > -\infty,\\
X & \text{ otherwise.}
\end{cases}
\end{equation}
Given a probability distribution $\mu$ and an initial configuration $X_0 \in \S$, the infinite-bin model $(X_n,\, n\geq 0)$ with distribution $\mu$ starting from $X_0 \in \S$ can be constructed through the following recursion equation, for $n\geq 0$
\begin{equation}\label{eqn:ibm}
X_{n+1}
:= \Phi_{\xi_{n+1}} (X_n),
\end{equation}
where $(\xi_n, n \geq 1)$ is an i.i.d. sequence of random variables with law $\mu$.

\begin{prop}\label{prop:bounded}
For any $A\in \bbN$, there exists a constant $C_A <\infty$ depending only on $A$ and the distribution $(\mu(n), n\geq 0)$ such that, for any initial configuration $X_0 \in \S$ satisfying
\[
\sup_{k \in \Z} X_0(k) \leq A,
\]
it holds that
\begin{equation}\label{eq:unifboundE}
\sup_{k \in \Z} \E(X_\infty(k)) < C_A.
\end{equation}
In particular, if $X_0 \in \S$ is a bounded configuration, then $X_\infty(k) < \infty$ a.s. for all $k\in \Z$ so every bin freezes.
\end{prop}

\begin{proof}
Fix a configuration $X_0 \in \S$ such that $\sup_{k \in \Z} X_0(k) \leq A$ and fix $k\in \Z$. We define $((T_n,r_n), n \geq 0)$ as the sequence of times at which a new ball is added to a bin strictly to the right of the bin $k-1$, together with the rank of its parent. Fixing $T_0
:= 0$, we set
\[
T_{n+1}
:= \inf\left\{m > T_n : \xi_m \leq \sum_{i=k-1}^\infty X_{m-1}(i) \right\} \quad \text{and} \quad r_n
:= \xi_{T_n}.
\]
Denoting $M_n$ the total number of balls in the bins strictly to the right of $k-1$ at time $T_n$, we have that $M_n
:= \sum_{i=k}^\infty X_{T_n}(i) = M_0 + n$, and that $r_n$ has law $\mu$ conditioned on being at most equal to $M_n + X_{T_{n}}(k-1)$.

For any $n_0 \in \bbN$, we have the upper bound
\begin{equation}\label{eq:upbound1}
X_{\infty}(k) = X_0(k) + \sum_{n=0}^{+\infty} \ind{r_n\,>\, M_n} \leq A + n_0 + \sum_{n=n_0}^{+\infty} \ind{r_n\,>\,M_n}.
\end{equation}
We recall the notation $\bar{\mu}(n) = \mu(\llb n,\infty\llbracketclose)$ for the tail of $\mu$. We fix $n_0$ large enough such that
\[
\alpha
:= \frac{\bar{\mu}(n_0)}{\mu(\llb 1,n_0\rrb)} < 1.
\]
Taking the expectation in~\eqref{eq:upbound1} conditionally on $\clH = \sigma(X_{T_n}(k-1), n \geq 0)$, we find that
\begin{align*}
\E(X_\infty(k) \mid \clH) & \leq A + n_0 + \sum_{n=n_0}^{+\infty} \P(r_n > M_n \mid \clH)\\
& = A + n_0 + \sum_{n=n_0}^{+\infty} \frac{\mu(\llb M_n+1,M_n+X_{T_n}(k-1)\rrb)}{\mu(\llb 1,M_n+X_{T_n}(k-1)\rrb)}\\
& \leq A + n_0 + \frac{\sum_{n=n_0}^\infty \mu(\llb M_0 + n + 1,M_0 + n + X_{\infty}(k-1)\rrb)}{\mu(\llb 1,n_0 \rrb)} \\
& \leq A + n_0 + \frac{\bar{\mu}(M_0 + n_0 + 1) X_\infty(k-1)}{\mu(\llb 1,n_0 \rrb)} \\
& \leq A + n_0 + \alpha X_\infty(k-1).
\end{align*}
Taking the expectation on both sides of the inequality above, we obtain a recursion inequality valid for all $k\in \Z$,
\begin{equation}\label{eq:recbounded}
\E(X_\infty(k)) \leq A + n_0 + \alpha \E(X_\infty(k-1)).
\end{equation}
Let $\ell\leq0$ and consider now the truncated configuration $X^{(\ell)}_0 = X_0 \ind{\cdot\,\geq\,\ell}$ where all the balls in the bins on the left of bin $\ell$ have been removed. We take $\ell$ to be small enough so that $X^{(\ell)}_0$ is non-empty. We denote by $X^{(\ell)}$ the IBM process starting from $X^{(\ell)}_0$, using the same sequence $(\xi_n)$ to construct $X^{(\ell)}_n$ recursively by~\eqref{eqn:ibm}. By definition of the dynamics, no ball can ever be added to bin $\ell$ so
\[
\E\left(X^{(\ell)}_\infty(\ell)\right) = X^{(\ell)}_0(\ell) \leq A.
\]
In view of~\eqref{eq:recbounded}, it follows by induction that, for all $k\geq \ell$, we have
\[
\E\left(X^{(\ell)}_\infty(k)\right) \leq C_A
:= \frac{A + n_0}{1-\alpha}.
\]
Finally, for each fixed $k$ and $n$, we have $\lim_{\ell \to -\infty} X_n^{(\ell)}(k) = X_n(k)$ a.s. Indeed, we have either $\sum_{k \in \Z} X_0(k) < \infty$, then $X_0 = X_0^\ell$ for $\ell$ large enough, or $\sum_{k \in \Z} X_0(k) = \infty$, then a.s., for $-\ell$ large enough we have $\max_{j \leq n} \xi_j \leq \sum_{k = \ell}^\infty X_0(k)$, therefore $X^{(\ell)}_n(k) = X_n(k)$. As a result, by Fatou's Lemma, for any $k\in\Z$,
\[
\E(X_n(k)) \leq \liminf_{l \to -\infty} \E\left(X_n^{(\ell)}(k)\right) \leq C_A.
\]
By monotone convergence, we conclude that, for any $k\in\Z$,
\[
\E(X_\infty(k)) = \lim_{n \to +\infty} \E(X_n(k)) \leq C_A.
\]
This completes the proof of~\eqref{eq:unifboundE} and also of Theorem~\ref{thm:unifbounded} since a random variable with finite expectation is finite a.s.
\end{proof}

\begin{thebibliography}{ČS24}
\bibitem[FK03]{FK} Sergey Foss and Takis Konstantopoulos, Extended renovation theory and limit theorems for stochastic ordered graphs, Markov Process. Relat. Fields 9 (2003), no. 3, 413–468.
\bibitem[KS64]{KochenStone64} Simon Kochen and Charles Stone, A note on the Borel–Cantelli lemma, Ill. J. Math. 8 (1964), 248–251.
\end{thebibliography}
\end{document}
